\documentclass[12pt]{article}
\usepackage[T1]{fontenc}
\usepackage[utf8]{inputenc}
\usepackage{lmodern}
\usepackage{textcomp}
\usepackage{enumitem}
\usepackage{geometry}
\geometry{margin=1in}
\begin{document}
\begin{center}
Answer Key - Mathematics - K-12 Level Assessment 17 \
\small (Answers are ordered exactly as in the question paper.)
\end{center}

\section*{Multiple Choice}
\begin{enumerate}[label=\arabic*.]
\item \textbf{Answer:} B
\\ \textit{Options:} A) 5; B) 10; C) 2; D) 12
\\ \textit{Note:} When simplifying $(x\sqrt{x^3})^4$, we first rewrite $\sqrt{x^3}$ as $x^{3/2}$, resulting in $(x \cdot x^{3/2})^4 = (x^{1 + 3/2})^4 = (x^{5/2})^4$, which simplifies to $x^{10}$.
\item \textbf{Answer:} D
\\ \textit{Options:} A) 0; B) 1; C) 2; D) 3
\\ \textit{Note:} The expression simplifies to 2 when \( a = 42 \), and since \( a^{2+1-3} = a^0 = 1 \), adding these gives \( 2 + 1 = 3 \).
\item \textbf{Answer:} C
\\ \textit{Options:} A) 252; B) 4680; C) 98770; D) 101170
\\ \textit{Note:} The correct answer, 98770, is obtained by using the formula for combinations, specifically $\dbinom{n}{k} = \frac{n!}{k!(n-k)!}$, where $\dbinom{85}{82}$ simplifies to $\dbinom{85}{3}$, which calculates to 98770.
\item \textbf{Answer:} C
\\ \textit{Options:} A) \ensuremath{\frac{125}{2916}}; B) \ensuremath{\frac{5}{16}}; C) \ensuremath{\frac{5}{54}}; D) \ensuremath{\frac{5}{55}}
\\ \textit{Note:} The correct answer is \(\frac{5}{54}\) because there are 3 prime numbers (2, 3, 5) and 3 composite numbers (4, 6) on a 6-sided die, and the probability of rolling exactly three primes and three composites can be calculated using combinations and the probabilities of each outcome.
\item \textbf{Answer:} D
\\ \textit{Options:} A) 458; B) 463; C) 547; D) 559
\\ \textit{Note:} The number 559 cannot be the original number before rounding because it rounds to 600, while numbers between 450 and 549 round to 500.
\end{enumerate}

\section*{True / False}
\begin{enumerate}[label=\arabic*.]
\setcounter{enumi}{5}
\item \textbf{Answer:} False
\\ \textit{Options:} A) True; B) False
\\ \textit{Note:} The number 179,912 in word form is correctly written as one hundred seventy-nine thousand nine hundred twelve, not one hundred seventy-nine thousand twelve.
\item \textbf{Answer:} False
\\ \textit{Options:} A) True; B) False
\\ \textit{Note:} The statement is false because the correct number that makes the equation 2/9 equal to 14/? is 63, not 15, as determined by cross-multiplying to find the unknown denominator.
\item \textbf{Answer:} False
\\ \textit{Options:} A) True; B) False
\\ \textit{Note:} The statement is false because the expression simplifies to $(3 \cdot 2^{20} + 2^{21}) \div 2^{17}$, which equals $10$, but the calculation is incorrect because it does not take into account proper simplification and evaluation of powers of 2.
\item \textbf{Answer:} False
\\ \textit{Options:} A) True; B) False
\\ \textit{Note:} The expression -\ensuremath{\dfrac{1}{-3}}\ensuremath{\cdot}\ensuremath{\cfrac{1}}\{\textasciitilde{}\ensuremath{\frac{1}{-3}}\textasciitilde{}\} simplifies to -3, not 1, because multiplying a negative fraction by the reciprocal of that fraction results in a negative value.
\item \textbf{Answer:} True
\\ \textit{Options:} A) True; B) False
\\ \textit{Note:} The statement is true because if 1 inch on the map represents 25 miles, then 5 inches would represent 5 times 25 miles, which equals 125 miles.
\end{enumerate}

\section*{Long Form Response}
\begin{enumerate}[label=\arabic*.]
\setcounter{enumi}{10}
\item \textbf{Answer:} Note that the numbers increase by 3 each time. Hence we increase by 3 a total of $\frac{32 - (-4)}{3} = 12$ times. But then there must be $12 + 1 = \boxed{13}$ numbers, since we need to also include the first number on the list.
\\ \textit{Note:} correct because it determines the total number of terms in the arithmetic sequence by calculating the number of increments of 3 from the starting value of -4 to the endpoint of 32, and then adding one to account for the initial term.
\item \textbf{Answer:} The concept of vertical axis of symmetry in vowels refers to the idea that certain vowel shapes can be divided into two mirror-image halves along a vertical line. Most vowels, such as A, O, and U, exhibit this symmetry because they look the same on both sides of the vertical line. However, the vowel E lacks this vertical axis of symmetry, as its shape is not identical when split down the middle. For instance, if you draw a vertical line through the center of E, the left and right sides do not match, demonstrating the absence of symmetry. Therefore, E is the only vowel that does not possess this characteristic.
\\ \textit{Note:} correct because it accurately identifies the vertical axis of symmetry in vowels, demonstrates that vowels like A, O, and U possess this symmetry, and clearly explains that the vowel E does not have this symmetry due to its asymmetrical shape when divided vertically.
\end{enumerate}

\vfill
\noindent\textit{End of Answer Key}
\end{document}